% GENERATED FILE. Edit canonical lessons, then run npm run generate:books.
% canonical-source-hash: fnv1a64:de0e27ca19992891
% canonical-lessons: SA-C528-ayojayati, SA-C528-samagri, SA-C528-sambharah, SA-C528-upakalpayati, SA-C528-pancangam

\chapter{Making a Plan}
\label{ch:sa-making-a-plan}
\hlchaptermodality{528}

\begin{chapteropening}
\textbf{By the end of this chapter:} \emph{I can say what I will organise if things work out.}

Complete the last lesson of chapter 528: I can say what I will organise if things work out.
\end{chapteropening}

\section[\texorpdfstring{\sk{आयोजयति} (āyojayati)}{आयोजयति (āyojayati)}]{\sk{आयोजयति} (āyojayati) — organises; arranges}

\label{lesson:SA-C528-ayojayati}



\begin{quote}
Before the new one: say the Sanskrit for cheers on, then the Sanskrit for a bowstring.
\end{quote}

\subsection*{You'll want to know: \sk{आयोजयति}}
\textbf{\sk{आयोजयति}} — \emph{āyojayati} — \textquotedblleft{}organises; arranges\textquotedblright{}.

\begin{quote}
\sk{यदि} \sk{भवान्} \sk{इच्छति}, \sk{तर्हि} \sk{अहं} \sk{सभाम्} \sk{आयोजयामि।}
\end{quote}

\begin{quote}
\emph{yadi bhavān icchati, tarhi ahaṃ sabhām āyojayāmi.}
\end{quote}

\begin{quote}
If you wish, I will organise a meeting.
\end{quote}

Past (he …): \sk{सः} \sk{उत्सवम्} \sk{आयोजितवान्}, \emph{saḥ utsavam āyojitavān}.

Future: \sk{आयोजयिष्यति}, \emph{āyojayiṣyati}.

\begin{cousinweb}
\textbf{\sk{आयोजयति}} (\emph{āyojayati}) is built on \emph{ā} and the causative of the root √yuj, \textquotedblleft{}to join\textquotedblright{}, which also gives \emph{yojayati}; MW gives \emph{āyojayati} as to join together, to form; \textquotedblleft{}to organise\textquotedblright{} is its modern use (MW).
\end{cousinweb}

\begin{grammarlens}[title={What you've built}]
One more word for saying what happens if.
\end{grammarlens}

\subsection*{Your turn}
\begin{itemize}
\raggedright
  \item \emph{Say it:} \emph{āyojayati}
  \item \emph{Say it:} \emph{āyojayati}, once more
  \item \emph{Say it:} \emph{yadi bhavān icchati, tarhi ahaṃ sabhām āyojayāmi.} — If you wish, I will organise a meeting.
  \item \emph{Recall:} say the Sanskrit for cheers on, then the Sanskrit for a bowstring, then say \emph{āyojayati} again
\end{itemize}

\begin{tcolorbox}[breakable,colback=teal!4,colframe=teal!35!black,title={Before you move on}]
What does \sk{आयोजयति} mean? (\textquotedblleft{}Organises; arranges\textquotedblright{}.) Say it once more.
\end{tcolorbox}
\section[\texorpdfstring{\sk{सामग्री} (sāmagrī)}{सामग्री (sāmagrī)}]{\sk{सामग्री} (sāmagrī) — materials}

\label{lesson:SA-C528-samagri}



\begin{quote}
Before the new one: say the Sanskrit for a bowstring, then the Sanskrit for organises.
\end{quote}

\subsection*{You'll want to know: \sk{सामग्री}}
\textbf{\sk{सामग्री}} — \emph{sāmagrī} — \textquotedblleft{}materials\textquotedblright{}.

\begin{quote}
\sk{यदि} \sk{सामग्री} \sk{सज्जा} \sk{अस्ति}, \sk{तर्हि} \sk{अहं} \sk{पायसं} \sk{पचामि।}
\end{quote}

\begin{quote}
\emph{yadi sāmagrī sajjā asti, tarhi ahaṃ pāyasaṃ pacāmi.}
\end{quote}

\begin{quote}
If the materials are ready, I'll make rice pudding.
\end{quote}

A feminine noun.

\begin{cousinweb}
\textbf{\sk{सामग्री}} (\emph{sāmagrī}) comes from \emph{samagra}, \textquotedblleft{}whole, complete\textquotedblright{}, as MW derives it: entirety, and so a complete set of implements or materials (MW).
\end{cousinweb}

\begin{grammarlens}[title={What you've built}]
One more word for saying what happens if.
\end{grammarlens}

\subsection*{Your turn}
\begin{itemize}
\raggedright
  \item \emph{Say it:} \emph{sāmagrī}
  \item \emph{Say it:} \emph{sāmagrī}, once more
  \item \emph{Say it:} \emph{yadi sāmagrī sajjā asti, tarhi ahaṃ pāyasaṃ pacāmi.} — If the materials are ready, I'll make rice pudding.
  \item \emph{Recall:} say the Sanskrit for a bowstring, then the Sanskrit for organises, then say \emph{sāmagrī} again
\end{itemize}

\begin{tcolorbox}[breakable,colback=teal!4,colframe=teal!35!black,title={Before you move on}]
What does \sk{सामग्री} mean? (\textquotedblleft{}Materials\textquotedblright{}.) Say it once more.
\end{tcolorbox}
\section[\texorpdfstring{\sk{सम्भारः} (sambhāraḥ)}{सम्भारः (sambhāraḥ)}]{\sk{सम्भारः} (sambhāraḥ) — supplies; provisions}

\label{lesson:SA-C528-sambharah}



\begin{quote}
Before the new one: say the Sanskrit for organises, then the Sanskrit for materials.
\end{quote}

\subsection*{You'll want to know: \sk{सम्भारः}}
\textbf{\sk{सम्भारः}} — \emph{sambhāraḥ} — \textquotedblleft{}supplies; provisions\textquotedblright{}.

\begin{quote}
\sk{यदि} \sk{अतिथयः} \sk{आगच्छन्ति}, \sk{तर्हि} \sk{सम्भारः} \sk{आवश्यकः} \sk{अस्ति।}
\end{quote}

\begin{quote}
\emph{yadi atithayaḥ āgacchanti, tarhi sambhāraḥ āvaśyakaḥ asti.}
\end{quote}

\begin{quote}
If guests are coming, we need supplies.
\end{quote}

A masculine noun.

\begin{cousinweb}
\textbf{\sk{सम्भारः}} (\emph{sambhāraḥ}) is built on \emph{sam}, \textquotedblleft{}together\textquotedblright{}, and the root √bhṛ, \textquotedblleft{}to bear, to bring\textquotedblright{}, which also gives \emph{bharati}; MW gives \emph{sambhāra} as preparation, provision, the materials needed (MW).
\end{cousinweb}

\begin{grammarlens}[title={What you've built}]
One more word for saying what happens if.
\end{grammarlens}

\subsection*{Your turn}
\begin{itemize}
\raggedright
  \item \emph{Say it:} \emph{sambhāraḥ}
  \item \emph{Say it:} \emph{sambhāraḥ}, once more
  \item \emph{Say it:} \emph{yadi atithayaḥ āgacchanti, tarhi sambhāraḥ āvaśyakaḥ asti.} — If guests are coming, we need supplies.
  \item \emph{Recall:} say the Sanskrit for organises, then the Sanskrit for materials, then say \emph{sambhāraḥ} again
\end{itemize}

\begin{tcolorbox}[breakable,colback=teal!4,colframe=teal!35!black,title={Before you move on}]
What does \sk{सम्भारः} mean? (\textquotedblleft{}Supplies; provisions\textquotedblright{}.) Say it once more.
\end{tcolorbox}
\section[\texorpdfstring{\sk{उपकल्पयति} (upakalpayati)}{उपकल्पयति (upakalpayati)}]{\sk{उपकल्पयति} (upakalpayati) — makes ready; prepares}

\label{lesson:SA-C528-upakalpayati}



\begin{quote}
Before the new one: say the Sanskrit for materials, then the Sanskrit for supplies.
\end{quote}

\subsection*{You'll want to know: \sk{उपकल्पयति}}
\textbf{\sk{उपकल्पयति}} — \emph{upakalpayati} — \textquotedblleft{}makes ready; prepares\textquotedblright{}.

\begin{quote}
\sk{यदि} \sk{श्वः} \sk{यात्रा} \sk{अस्ति}, \sk{तर्हि} \sk{अद्य} \sk{स्यूतम्} \sk{उपकल्पयतु।}
\end{quote}

\begin{quote}
\emph{yadi śvaḥ yātrā asti, tarhi adya syūtam upakalpayatu.}
\end{quote}

\begin{quote}
If the trip is tomorrow, get your bag ready today.
\end{quote}

Past (he …): \sk{सः} \sk{स्यूतम्} \sk{उपकल्पितवान्}, \emph{saḥ syūtam upakalpitavān}.

\begin{cousinweb}
\textbf{\sk{उपकल्पयति}} (\emph{upakalpayati}) is built on \emph{upa}, \textquotedblleft{}near\textquotedblright{}, and the causative of the root √kḷp, \textquotedblleft{}to be fit, to be in order\textquotedblright{}; MW gives \emph{upakalpayati} as to prepare, to make ready, to equip (MW).
\end{cousinweb}

\begin{grammarlens}[title={What you've built}]
One more word for saying what happens if.
\end{grammarlens}

\subsection*{Your turn}
\begin{itemize}
\raggedright
  \item \emph{Say it:} \emph{upakalpayati}
  \item \emph{Say it:} \emph{upakalpayati}, once more
  \item \emph{Say it:} \emph{yadi śvaḥ yātrā asti, tarhi adya syūtam upakalpayatu.} — If the trip is tomorrow, get your bag ready today.
  \item \emph{Recall:} say the Sanskrit for materials, then the Sanskrit for supplies, then say \emph{upakalpayati} again
\end{itemize}

\begin{tcolorbox}[breakable,colback=teal!4,colframe=teal!35!black,title={Before you move on}]
What does \sk{उपकल्पयति} mean? (\textquotedblleft{}Makes ready; prepares\textquotedblright{}.) Say it once more.
\end{tcolorbox}
\section[\texorpdfstring{\sk{पञ्चांगम्} (pañcāṅgam)}{पञ्चांगम् (pañcāṅgam)}]{\sk{पञ्चांगम्} (pañcāṅgam) — an almanac}

\label{lesson:SA-C528-pancangam}



\begin{quote}
Before the new one: say the Sanskrit for supplies, then the Sanskrit for makes ready.
\end{quote}

\subsection*{You'll want to know: \sk{पञ्चांगम्}}
\textbf{\sk{पञ्चांगम्}} — \emph{pañcāṅgam} — \textquotedblleft{}an almanac\textquotedblright{}.

\begin{quote}
\sk{यदि} \sk{पञ्चांगं} \sk{शुभं} \sk{दिनं} \sk{वदति}, \sk{तर्हि} \sk{विवाहः} \sk{भविष्यति।}
\end{quote}

\begin{quote}
\emph{yadi pañcāṅgaṃ śubhaṃ dinaṃ vadati, tarhi vivāhaḥ bhaviṣyati.}
\end{quote}

\begin{quote}
If the almanac gives a lucky day, the wedding will take place.
\end{quote}

A neuter noun.

\begin{cousinweb}
\textbf{\sk{पञ्चांगम्}} (\emph{pañcāṅgam}) is built from \emph{pañca}, \textquotedblleft{}five\textquotedblright{}, and \emph{aṅga}, \textquotedblleft{}a part\textquotedblright{}; MW gives \emph{pañcāṅga} as a calendar or almanac, which treats of five things, among them the days and the lunar days (MW).
\end{cousinweb}

\begin{grammarlens}[title={What you've built}]
One more word for saying what happens if.
\end{grammarlens}

\subsection*{Your turn}
\begin{itemize}
\raggedright
  \item \emph{Say it:} \emph{pañcāṅgam}
  \item \emph{Say it:} \emph{pañcāṅgam}, once more
  \item \emph{Say it:} \emph{yadi pañcāṅgaṃ śubhaṃ dinaṃ vadati, tarhi vivāhaḥ bhaviṣyati.} — If the almanac gives a lucky day, the wedding will take place.
  \item \emph{Recall:} say the Sanskrit for supplies, then the Sanskrit for makes ready, then say \emph{pañcāṅgam} again
\end{itemize}

\begin{tcolorbox}[breakable,colback=teal!4,colframe=teal!35!black,title={Before you move on}]
What does \sk{पञ्चांगम्} mean? (\textquotedblleft{}An almanac\textquotedblright{}.) Say it once more.
\end{tcolorbox}
